\subsection{The central cusp defect and the open descent image}

Retain $U,U^{\mathrm c},M$ of (B.S1) for the particular
six-coordinate quotient.  Thus $M$ is finite by B.2.15.

\paragraph{The actual cusp defect.}
Let $\alpha$ be a positive $\tau$-fixed automorphism of $U$,
with pure frames $\beta,A\in J$.  Choose any pure lifts
$\widetilde\beta,\widetilde A$ to $U^{\mathrm c}$.
The assignments
\[
 \widetilde\alpha(s)=\widetilde\beta s\widetilde\beta^{-1},
 \qquad
 \widetilde\alpha(r)=\widetilde A r\widetilde A^{-1}            \tag{B.S9}
\]
define an actual automorphism of $U^{\mathrm c}$ above
$\alpha$.  Indeed they define a homomorphism from $B$,
send $K$ into central $M$, and therefore kill
$\overline{[B,K]}$.  Let $V$ be its closed image.  Since
the induced map on $U$ is onto, $U^{\mathrm c}=VM$.
The kernel $M$ is central and lies in
$(U^{\mathrm c})'$, because $K\subseteq B'$.
Consequently $(U^{\mathrm c})'=V'\subseteq V$, whence
$M\subseteq V$ and $V=U^{\mathrm c}$.  Hopficity of
the finitely generated profinite group $U^{\mathrm c}$
makes the endomorphism an automorphism.  Its induced
automorphism of $U$ shows that it preserves $M$.

Its action on the named cusp is exactly
\[
 \widetilde\alpha(\tau)=\tau d,\qquad
 d=\tau^{-1}\widetilde\beta s\widetilde\beta^{-1}
                       \widetilde A r\widetilde A^{-1}\in M.
                                                                  \tag{B.S10}
\]
Changing a chosen lift of either fixed $U$-frame changes
it by a central element of $M$ and hence leaves (B.S9)
and $d$ unchanged.  The multiplier-one cusp condition is
$d=1$; it is not supplied by positivity of the two torsion
images alone.

There are exact continuous centralizer maps
\[
\begin{aligned}
 \kappa_s:C_J(s)&\longrightarrow M[2],&
          c&\longmapsto\widetilde c s\widetilde c^{-1}s^{-1},\\
 \kappa_r:C_J(r)&\longrightarrow M[3],&
          c&\longmapsto\widetilde c r\widetilde c^{-1}r^{-1}.
\end{aligned}                                                   \tag{B.S11}
\]
Their values are independent of the lifts, since $M$ is
central.  The commutators are central, so the maps are
homomorphisms; the source torsion relation makes their
values have order dividing two or three respectively.
Every other pure positive frame of the same $\alpha$ is
$\beta c_s$ or $Ac_r$ with the indicated centralizer
element.  Its defect is exactly
$d\kappa_s(c_s)\kappa_r(c_r)$.  Thus the full set of defects
for this fixed $\alpha$ is the literal coset
$d\,\operatorname{im}\kappa_s\operatorname{im}\kappa_r$ in $M$.
This identifies the remaining finite cusp obstruction
without dropping the named $\tau$.

\begin{result}{B.2.16}{Proved}{The descent endomorphism has open image}
The subgroup $D_P(H)$ is open, and hence of finite index,
in $H$.  The image of $D_P$ on the complete gentle group
$\widehat{GT}_0$ is likewise open and of finite index.
\end{result}

\begin{proof}
The unique actual pure torsion frames $\beta_g,A_g$ of a
normalized $g\in H$ depend continuously on $g$.  Their
uniqueness follows from the actual torsion centralizers
$C_B(s)=\langle s\rangle$ and
$C_B(r)=\langle r\rangle$, which have trivial intersection
with $F$.  Continuity also follows directly from their
closed graphs: the frame equations define a compact
subset of $H\times F$, and its projection to $H$ is a
continuous bijection onto a Hausdorff space.

Take the open identity neighborhood $W\subseteq H$ on
which both target frames reduce to one in $\mathcal T$.
For these targets choose $a=l=0$.  The tuple (B.R12)
then reduces to the identity in $\mathcal T^6$ and lies
in $J'$ by (B.R4).  This is precisely the successful
exponent in the identity row of B.2.12.  The construction
of B.2.13 therefore gives actual automorphisms
$\alpha_U(g)$ depending continuously on $g\in W$;
their values on the two source generators are the
continuous word formulas (B.R15).  Each acts trivially
on $E=U/L_0^6$.

The pure positive frames for $\alpha_U(g)$ can also be
chosen continuously on this neighborhood.  Use the
identity finite frame in $E$, and use the explicit
coordinate solution of B.2.13 in $L_0^6$: on each
$s$-pair or $r$-triple set the first coordinate to one
and solve the remaining coordinates successively.
All factor isomorphisms in this recursion are the
fixed source conjugation maps.  The solution is thus
a continuous group-word function of the relevant
coordinates of $\alpha_U(g)(s)s^{-1}$ or
$\alpha_U(g)(r)r^{-1}$.  It equals one at $g=1$.
Denote these actual frames by $\beta_U(g),A_U(g)$.

No continuous section of $U^{\mathrm c}\to U$ is needed
to form (B.S9).  Centrality of $M$ makes conjugation
by a lift of an element of $J$ independent of that
lift.  The continuous conjugation map on the pure
subgroup of $U^{\mathrm c}$ therefore factors through
the quotient $J$.  Thus the two elements in (B.S9),
and consequently the defect $d(g)$ in (B.S10), are
continuous functions of the chosen $U$-frames.
They give a continuous map $d:W\to M$ with $d(1)=1$.
Since $M$ is finite, $W_0=d^{-1}(1)$ is an open
identity neighborhood.

For every $g\in W_0$, (B.S9) is a positive
$\tau$-fixed automorphism of $U^{\mathrm c}$ preserving
$M$.  The exact image equality B.2.14 supplies a full
positive $\tau$-fixed automorphism $\Phi$ of $B$ with
the prescribed $U$-action.  Its literal descendant
through $\varphi_P$ is
$\operatorname{Inn}(\beta_g^{-1})\alpha_g$, by the
first coordinate of (B.R19).  Finally apply the exact
commutator normalization (B.M13): replacing $\Phi$
by $\operatorname{Inn}(x^b)\Phi$ gives an element
$h\in H$ and changes the descendant only by
$\operatorname{Inn}(\varphi_P(x)^b)
=\operatorname{Inn}(s^b)$.  Its outer target class is
still $g$.  Faithfulness of the normalized target
parameter therefore gives $D_P(h)=g$.

This proves $W_0\subseteq D_P(H)$.  A subgroup
containing an open identity neighborhood is open;
compactness of $H$ makes its index finite.  For the
complete group, equation (B.Q9) gives
$\operatorname{im}D_P=G_{\mathbb Q}D_P(H)$ and its
intersection with $H$ is $D_P(H)$.  Arithmetic
cyclotomic surjectivity then identifies its index
in $\widehat{GT}_0$ with the finite index
$[H:D_P(H)]$.  The full image is closed, being the
continuous image of a compact group.  A closed
finite-index subgroup of a profinite group is open,
proving the second assertion.
\end{proof}

The finite index in this result has not been computed.
In particular, the proof leaves the index-one question
and independence of all later descent characters open.
